\documentclass[11pt]{article}

% --------------------------------------------------
% Packages
% --------------------------------------------------
\usepackage[margin=1in]{geometry}
\usepackage{amsmath, amssymb, amsthm}
\usepackage{mathtools}
\usepackage{graphicx}
\usepackage{float}
\usepackage{enumitem}
\usepackage{hyperref}
\usepackage{xcolor}
\usepackage{tikz}
\usetikzlibrary{arrows.meta, positioning, calc}

% Pseudocode and algorithms
\usepackage{algorithm}
\usepackage{algpseudocode}

% Better fonts (optional)
\usepackage{lmodern}

% --------------------------------------------------
% Hyperref setup
% --------------------------------------------------
\hypersetup{
    colorlinks=true,
    linkcolor=blue,
    citecolor=blue,
    urlcolor=blue
}

% --------------------------------------------------
% Theorem-like environments
% --------------------------------------------------
\theoremstyle{definition}
\newtheorem{definition}{Definition}[section]
\newtheorem{theorem}{Theorem}[section]
\newtheorem{claim}{Claim}[section]
\newtheorem{remark}{Remark}[section]
\newtheorem{recall}{Recall}[section]

\theoremstyle{proof}
% Proof environment is built-in via amsthm

% --------------------------------------------------
% Custom commands
% --------------------------------------------------
\newcommand{\Gen}{\mathsf{Gen}}
\newcommand{\Enc}{\mathsf{Enc}}
\newcommand{\Dec}{\mathsf{Dec}}
\newcommand{\Sig}{\mathsf{Sig}}
\newcommand{\Ver}{\mathsf{Ver}}
\newcommand{\negl}{\mathsf{negl}}

% --------------------------------------------------
% Document
% --------------------------------------------------
\begin{document}

\title{Public Key Crypto notes - 12}
\author{By contributors}
\date{May 04. 2026}
\maketitle

\section{Elliptic Curve ElGamal Signature Scheme}

Let $p$ be a prime, $E$ be an elliptic curve over $\mathbb{Z}_p$, $P \in E$ of prime order $l$. Also, let $f: E \rightarrow \mathbb{Z}$ (typical choice: $f(x, y) = x$).

\begin{itemize}
	\item KeyGen:
	\begin{itemize}
		\item $sk = a \in_R \mathbb{Z}_l$, $pk = (p, E, P, l, a \cdot P = Q)$.
	\end{itemize}
	\item Sign:
	\begin{itemize}
		\item Let $m \in \mathbb{Z}_l$,
		\item Choose $r \in_R \mathbb{Z}_l, R = rP \in E$,
		\item Compute $s = k^{-1}(m - a \cdot f(R))$,
		\item Signature: $(m, R, s)$
	\end{itemize}
	\item Verify:
	\begin{itemize}
		\item Let $pk = Q, sign = (m, R, S)$
		\item $v_1 = sR + f(R) \cdot Q$
		\item $v_2 = m \cdot P$
		\item Verify: $v_1 \stackrel{?}{=} v_2$
	\end{itemize}
	\item Correctness:
	\begin{align}
		v_1 &= sR + f(R)Q = k^{-1}(m-af(R)) kP + f(R)aP =\\
		&= mP-af(R)P+af(R)P = mP = v_2
	\end{align}
\end{itemize}

\section{Elliptic Curve Digital Signature Algorithm (ECDSA)}

\begin{itemize}
	\item Standardized by NIST.
	\item KeyGen:
	\begin{itemize}
		\item $sk = a \in_R \mathbb{Z}_l, pk = Q = aP \in E$
	\end{itemize}
	\item Sign:
	\begin{itemize}
		\item Let $m \in \mathbb{Z}_l$ be the message.
		\item Choose $k \in_R \mathbb{Z}_l, R = kP \in E$.
		\item Compute $s = k^{-1} (m + af(R))$
		\item Signature: $(m, R, s)$
	\end{itemize}
	\item Verify:
	\begin{itemize}
		\item $u_1 = s^{-1} \cdot m \in \mathbb{Z}_l$
		\item $u_2 = s^{-1} \cdot f(R) \in \mathbb{Z}_l$
		\item $v = u_1 P + u_2 Q$
	\end{itemize}
	\item Correctness:
	\begin{align}
		v &= u_1 P + u_2 Q \\
		&= s^{-1} \cdot mP + s^{-1} \cdot f(R)a \cdot P \\
		&= (s^{-1} \cdot m + s^{-1} \cdot f(R)a) P \\
		&= k (m + af(R))^{-1} (m + a f(R)) \cdot P \\
		&= kP = R
	\end{align}
\end{itemize}

\begin{remark}[Computational costs of Verify]
	\leavevmode

	EC ElGamal: 3 scalar multiplications, 1 addition.
	
	ECDSA: 2 scalar multiplications, 1 addition on $E$.
	
	ECDSA is much faster.
\end{remark}

\begin{remark}
	One can consider elliptic curves over finite field $\mathbb{F}_q$ ($q$ is a prime power) instead of $\mathbb{Z}_p$.
	
	However, if $q = 2^n$ or $q = 3^n$, we need a different curve equation.
\end{remark}

\section{Pairing-based crypto}

Let $E$ be an elliptic curve over $\mathbb{Z}_p$, and for $n | |E|$, let $E[n] = \{P \in E, nP = \infty\}$.

If $n$ is large enough, then $E[n]$ can be almost as large as $E$, e.g. $|E| = l$ is a prime $E[l] = E$.

\begin{definition}[Pairing]
	A pairing is a map $e: E[n] \times E[n] \rightarrow \mathbb{Z}_p$ with:
	
	\begin{enumerate}
		\item $e$ is bilinear. $e(aP, Q) = e(P, aQ) = e(P, Q)^a$
		\item $e$ is non-trivial. If $e(P, Q) = 1$ for all $Q$ or $e(Q, P) = P$ for all $Q \rightarrow P = \infty$ 
	\end{enumerate}
\end{definition}

\begin{remark}
	For appropriate choice of parameters (e.g. $n | p-1$), there are curves $E$ over $\mathbb{Z}_p$ such that a pairing exists.
\end{remark}

Having pairing, one can perform cryptographic attacks.

\begin{enumerate}
	\item If there is a pairing, then the Decisional Diffie-Hellman problem is easy.
	\begin{itemize}
		\item Given $aP, bP, Q \rightarrow$ decide whether $abP \stackrel{?}{=}Q$
		\item Idea: $e(aP, bP) = e(P, abP) \stackrel{?}{=} e(P, Q)$
	\end{itemize}
	\item If one can solve DLog in $\mathbb{Z}_p$, then they can do it on $E$ too.
	\begin{itemize}
		\item Given $P, aP = Q \rightarrow a = ?$
		\item Idea: $e(P, Q) = e(P, aP) = e(P, P)^a$. If $g = e(P, P) \in \mathbb{Z}_p$ given $e(P, Q) = g^a \in \mathbb{Z}_p$.
		\item So $p$ has to be the same size as in classical (not elliptic curve) crypto, i.e. $p \approx 3000-4000$ bits.
	\end{itemize}
\end{enumerate}

\subsection{Protocol 1}

3-party Diffie-Hellman key exchange protocol.

\begin{itemize}
	\item Person $A$: $a \in_R \mathbb{Z}_n$, sends $s_AP$.
	\item Person $B$: $b \in_R \mathbb{Z}_n$, sends $s_BP$.
	\item Person $C$: $c \in_R \mathbb{Z}_n$, sends $s_CP$.
	\item Then $e(bP, cP)^a = e(aP, cP)^b, e(aP, bP)^c = e(P, P)^{abc}$.
\end{itemize}

\subsection{Protocol 2 - Identity-based encryption}

Goal: Design a PKE such that the public key is an identifier (ID).

\begin{definition}[Boneh-Franklin IBE]
	\begin{itemize}
		\item $A, B$ users
		\item TA: Trusted Authority
		\item Protocols:
		\begin{itemize}
			\item Setup,
			\item Extraction,
			\item Encryption,
			\item Decryption
		\end{itemize}
		\item Setup:
		\begin{itemize}
			\item TA chooses:
			\begin{itemize}
				\item Primes $p, l$, an elliptic curve $E$ over $\mathbb{Z}_p$ such that $\exists e: E[n] \times E[n] \rightarrow \mathbb{Z}_p$ and a point $P \in E[n]$ of order $l$.
				\item $H: \{0, 1\}^* \rightarrow P \in E[l]$
				\item $s \in_R \mathbb{Z}_l, P_{pub} = sP$, where $s$ is secret, $P_{pub}$ is public.
			\end{itemize}
		\end{itemize}
	\end{itemize}
\end{definition}

Extraction:

\begin{itemize}
	\item Let Bob's ID be $ID_B$.
	\item TA computes: $Q_B = H(ID_B), D_B = sQ_B$.
	\item Bob's $sk = D_B$
\end{itemize}

Encryption:

\begin{itemize}
	\item Let $ID_B$ be the ID of Bob.
	\item $A$ computes: $Q_B = H(ID_B), g_B = e(Q_B, P_{pub})$
	\item $r \in_R \mathbb{Z}_l$, cipher text: $c = (rP \in E, m \cdot g_B^r \in \mathbb{Z}_l)$ on message $m$.
\end{itemize}

Decryption:

\begin{itemize}
	\item $B$ is given $c = (U, v)$ having $D_B: h_B = e(D_B, U)$
	\item $m = \frac{v}{h_B}$
\end{itemize}

Correctness: 

\begin{align}
	h_B &= e(D_B, U) = \\
	&= e(sQ_B, rP) = \\
	&= r(Q_B, sP)^r = \\
	&= g_B^r
\end{align}

And:

$$
\frac{v}{h_B} = \frac{m \cdot g_B^r}{g_B^r} = m
$$

\end{document}